\section{Calling the Model from Python}
\label{sec:python}

Python is used to run the MiniZinc model on the benchmark instances. The script
loads the model and selects the solver:

\vspace{0.3cm}
\begin{lstlisting}[language=Python]
model    = Model("model.mzn")
solver   = Solver.lookup(solver_name)
instance = Instance(solver, model)
\end{lstlisting}
\vspace{0.3cm}

It then sets the instance parameters:

\vspace{0.3cm}
\begin{lstlisting}[language=Python]
instance["n_boxes"]       = pallet_info["n_boxes"]
instance["boxes_x_dim"]   = boxes_x_dim
instance["boxes_per_row"] = pallet_info["boxes_per_row"]
instance["pallet_rows"]   = pallet_info["pallet_rows"]
instance["grip_size"]     = grip_size
\end{lstlisting}
\vspace{0.3cm}

After the parameters are set, the script calls the solver:

\vspace{0.3cm}
\begin{lstlisting}[language=Python]
result = instance.solve()
\end{lstlisting}
\vspace{0.3cm}

The output arrays \texttt{chunk\_start}, \texttt{chunk\_len},
\texttt{plate}, \texttt{y\_row}, and \texttt{x\_start} are then read from the
\texttt{result} map. They give the grouping of the boxes, the target placement
of each group, and the release schedule. Since the belt is FIFO, the release
order is given directly by the chunk index.